\section{Introduction}

The Witten--Reshetikhin--Turaev invariant of a fixed closed three-manifold need not have a logarithm-free large-level expansion. We study a concrete uncolored family in which two opposite order-two caps are attached to every vertex of a tree plumbing. All cap colors are summed. The only final observable is the closed invariant. There are no Wilson insertions, boundary filters, or level-dependent manifolds.

The smallest nontrivial case is the six-component H plumbing of Figure~\ref{fig:plumbing}. Its sphere-normalized invariant grows as \((-1)^r r^2\log r/\pi^2\). Its exact finite reduction gives the entire logarithmic coefficient and the next constant, on both level parities. Attaching a pendant to each vertex of a fixed backbone tree produces arbitrarily high leading log powers. Fixed integer framings cut the polytope that determines the top coefficient; when that coefficient vanishes, the theorem gives an upper scale rather than a first nonzero term. A different exact transform applies to plain paths. The even and odd lengths have different dominant charge sectors, and odd paths have no leading logarithm. On an arbitrary zero-framed paired-cap tree, the exponent and the log degree are governed by the maximum matching and the maximizing face of a polytope defined by the edge inequalities.

Logarithmic terms already occur in the form of certain mapping-torus asymptotics \cite{AP}. Our obstruction concerns the particular all-closed finite-phase Laurent--Puiseux formulation in \cite[Conjecture~1.1]{AHLMS}, taken literally without logarithms. It does not concern that work's theorem for Seifert integral homology spheres. General exact Seifert formulas \cite{HT} provide another point of comparison with the elementary H-plumbing reduction. We make no priority claim about the appearance of logarithms or the earliest computation of any particular manifold. The proof uses finite modular matrices and actual plumbing surgery; it does not presume a semiclassical measure on singular flat moduli.

Section~\ref{sec:conventions} fixes the normalization. We prove the H calculation and then the framed corona result in the same conventions. Sections~\ref{sec:paths} and \ref{sec:trees} give the path transfer theorem and the general tree face law. We end with the scope of the obstruction and the integer-spin counterpart.
